% ID: SEL-D0007-BB-E03
% ESTADO: validada. S008, 2026-11-07; tres pasadas del mismo autor y revisión matemática, didáctica y visual.
\textbf{a) Monotonía y extremos.}\par\nopagebreak[4]
El dominio es $\mathbb R$ y $f(x)=x(x-2)^2$. La derivada es
\[
f'(x)=3x^2-8x+4=(3x-2)(x-2).
\]
Es positiva en $(-\infty,2/3)$ y $(2,\infty)$ y negativa en $(2/3,2)$. Así, $f$ crece en los dos intervalos de signo positivo y decrece entre 2/3 y 2. Los cambios de signo dan
\[
\boxed{\text{máximo relativo }(2/3,32/27),\qquad\text{mínimo relativo }(2,0).}
\]
No hay extremos absolutos sobre $\mathbb R$, porque la cúbica no está acotada.

\textbf{b) Representación.}\par\nopagebreak[4]
Sus cortes con el eje horizontal son $(0,0)$ y $(2,0)$; en 2 la raíz es doble y la curva toca el eje sin cruzarlo. Es negativa para $x<0$ y no negativa para $x\geq0$. Además, $f''(x)=6x-8$ cambia de signo en 4/3, con inflexión en $(4/3,16/27)$. En los extremos del eje real, $f(x)\to-\infty$ cuando $x\to-\infty$ y $f(x)\to\infty$ cuando $x\to\infty$.
\begin{center}\includegraphics[width=0.62\linewidth]{figuras/s008-cubica.pdf}\end{center}

\textbf{c) Integral indefinida.}\par\nopagebreak[4]
Por la regla de potencias y la linealidad,
\[
\boxed{\int f(x)\,dx=\frac{x^4}{4}-\frac43x^3+2x^2+K.}
\]
Derivando la expresión se recupera $x^3-4x^2+4x$.

\textbf{d) Área del recinto acotado.}\par\nopagebreak[4]
El único recinto acotado con el eje horizontal está entre las raíces 0 y 2. Como $f\geq0$ en $[0,2]$, el área geométrica coincide con la integral:
\[
S=\int_0^2 f(x)\,dx
=\left[\frac{x^4}{4}-\frac43x^3+2x^2\right]_0^2
=4-\frac{32}{3}+8=\boxed{\frac43\text{ unidades cuadradas}.}
\]
No se extiende el área hacia las ramas no acotadas de la curva.
